\documentclass[10pt,a4paper]{article}
\usepackage[utf8]{inputenc}
\usepackage[T1]{fontenc}
\usepackage[french]{babel}
\usepackage{lmodern}
\usepackage[left=1.5cm,right=1.5cm,top=1.25cm,bottom=1.25cm]{geometry}
\usepackage{xcolor}
\usepackage{hyperref}
\usepackage{enumitem}
\usepackage{parskip}
\usepackage{fontawesome5}
\usepackage{titlesec}
\usepackage{fancyhdr}
\usepackage{lastpage}

% Couleurs DevSecOps (Orange brûlé & Slate)
\definecolor{primary}{RGB}{194, 65, 12}       % Orange brûlé (#c2410c)
\definecolor{darkgray}{RGB}{30, 41, 59}       % Texte principal (#1e293b)
\definecolor{slateGray}{RGB}{71, 85, 105}     % Texte secondaire (#475569)
\definecolor{ruleGray}{RGB}{226, 232, 240}    % Ligne séparatrice

\hypersetup{
    colorlinks=true,
    linkcolor=primary,
    urlcolor=primary,
    pdfauthor={Astrid-Aurélien NKUMBE ENONGENE},
    pdftitle={CV Lead DevSecOps et Cloud Architect --- Astrid-Aurélien NKUMBE}
}

% Format des titres
\titleformat{\section}
  {\large\bfseries\color{primary}}
  {}{0em}
  {}[\vspace{-0.35em}\textcolor{primary}{\rule{\linewidth}{1.2pt}}\vspace{-0.25em}]

\titlespacing*{\section}{0pt}{0.85em}{0.4em}

% En-tête et pied de page
\pagestyle{fancy}
\fancyhf{}
\renewcommand{\headrulewidth}{0pt}
\renewcommand{\footrulewidth}{0.4pt}
\renewcommand{\footrule}{\hbox to\headwidth{\color{ruleGray}\leaders\hrule height \footrulewidth\hfill}}
\fancyfoot[L]{\footnotesize\color{slateGray}Astrid-Aurélien NKUMBE --- CV Lead DevSecOps \& Cloud}
\fancyfoot[C]{}
\fancyfoot[R]{\footnotesize\color{slateGray}Page \thepage\ / \pageref{LastPage}}

\color{darkgray}

\begin{document}

% =========================================================================
% EN-TÊTE PRINCIPAL
% =========================================================================
\begin{center}
    {\LARGE\bfseries\color{primary} Astrid-Aurélien NKUMBE ENONGENE}\\[0.25em]
    {\large\bfseries\color{darkgray} Lead DevSecOps Engineer $\cdot$ Cloud Infrastructure Architect}\\[0.4em]
    {\footnotesize\color{slateGray}
        \faMapMarker*\ Région Parisienne (Cergy / Paris) --- Mobile France \& Remote \quad$\cdot$\quad
        \faPhone*\ +33 7 44 58 45 62 \quad$\cdot$\quad
        \faEnvelope\ \href{mailto:nkumbeaurelien@hotmail.com}{\texttt{nkumbeaurelien@hotmail.com}}\\[0.25em]
        \faGlobe\ \href{https://nkaurelien.kamitbrains.fr}{\texttt{nkaurelien.kamitbrains.fr}} \quad$\cdot$\quad
        \faGithub\ \href{https://github.com/nkaurelien}{\texttt{github.com/nkaurelien}} \quad$\cdot$\quad
        \faLinkedin\ \href{https://www.linkedin.com/in/nkaurelien/}{\texttt{linkedin.com/in/nkaurelien}}
    }
\end{center}

\vspace{0.1em}
\textcolor{ruleGray}{\rule{\linewidth}{1pt}}
\vspace{0.35em}

% =========================================================================
% PROFIL
% =========================================================================
\section{\faLightbulb\ Profil Professionnel}
Ingénieur DevSecOps et Architecte Cloud avec \textbf{plus de 7 ans d'expérience (dont $\sim$5 ans en DevSecOps, Cloud et Platform Engineering)}, spécialisé dans l'automatisation intégrale des pipelines CI/CD, l'orchestration de conteneurs (Docker Swarm, Kubernetes, Argo CD) et la sécurité applicative (SAST/SCA Snyk, API Gateways, conformité HDS). Solide double compétence en développement logiciel (Python, Node.js, PHP) et en observabilité distribuée (SigNoz, Grafana, ELK Stack, Prometheus).

% =========================================================================
% COMPÉTENCES TECHNIQUES
% =========================================================================
\section{\faCogs\ Compétences DevSecOps \& Cloud}
\begin{itemize}[leftmargin=1.4em, itemsep=0.25em, topsep=0.15em]
    \item \textbf{Conteneurs \& Orchestration :} Docker, Docker Compose, Docker Swarm, Kubernetes (K8s), Argo CD, Portainer.
    \item \textbf{CI/CD \& Automatisation :} GitLab CI/CD, GitHub Actions, Azure DevOps, Ansible, Terraform, scripts Shell/Bash.
    \item \textbf{Sécurité Applicative (SecOps) :} Snyk (SAST/SCA), Gitleaks, KONG API Gateway, Nginx, Traefik, LDAP, Active Directory, Auth0, SSL/Certbot, RBAC.
    \item \textbf{Observabilité \& Métrologie :} SigNoz, Grafana, Loki, Alloy, Netdata, Prometheus, ELK Stack (Logstash, Elasticsearch, Kibana), Sentry.
    \item \textbf{Cloud Providers \& Plateformes :} AWS (EC2, S3, SQS, SNS), GCP (Cloud Run, GCS, Vertex AI), Azure (App Services, OpenAI), Supabase.
    \item \textbf{Systèmes \& Middleware :} Linux (Ubuntu, Debian), Software AG webMethods 10.x, Oracle 23ai, Postfix (SMTP Relay), PM2, Supervisor.
\end{itemize}

% =========================================================================
% EXPÉRIENCES PROFESSIONNELLES
% =========================================================================
\section{\faBriefcase\ Expériences DevSecOps \& Infrastructures}

\textbf{\large DATA2INNOV} \hfill {\bfseries\small\color{slateGray} Mai 2025 -- Présent} \\
\emph{\color{primary}Lead DevSecOps \& Fullstack Architect} \hfill {\small France}
\begin{itemize}[leftmargin=1.4em, itemsep=0.18em, topsep=0.15em]
    \item Conception de l'infrastructure de déploiement conteneurisée sécurisée avec Docker Swarm et Ansible.
    \item Pipeline CI/CD complet sous GitLab CI : exécution des tests, builds d'images multi-arch et push vers GitLab Container Registry.
    \item Déploiement de la métrologie et observabilité temps réel : SigNoz, Kafka et pipelines Logstash pour flux physiologiques IoT.
    \item Durcissement et sécurisation du réseau IoT (acquisition sécurisée BLE / Wi-Fi) et reverse proxies Nginx.
    \item Mise à jour du support et auteur des fiches pratiques (TP) et fichiers d'évaluation --- Formation \textbf{Cloud Computing \& Conteneurisation}.
\end{itemize}

\vspace{0.35em}

\textbf{\large YSO Conseils / KAZAC} \hfill {\bfseries\small\color{slateGray} Octobre 2023 -- Septembre 2024} \\
\emph{\color{primary}DevSecOps \& Cloud Infrastructure (SmartDataPay)} \hfill {\small France}
\begin{itemize}[leftmargin=1.4em, itemsep=0.18em, topsep=0.15em]
    \item Automatisation des déploiements d'applications web et IA sur le Cloud Azure avec GitLab CI.
    \item Déploiement d'environnements vectoriels d'embeddings sécurisés (Supabase pgvector, Databricks Notebooks, Docker).
    \item Automatisation de workflows SaaS avec Make/N8N et sécurisation des accès applicatifs via Auth0.
\end{itemize}

\newpage

\textbf{\large Korée Africa} \hfill {\bfseries\small\color{slateGray} Décembre 2021 -- Septembre 2023} \\
\emph{\color{primary}Lead DevSecOps \& Backend Architect} \hfill {\small Remote}
\begin{itemize}[leftmargin=1.4em, itemsep=0.18em, topsep=0.15em]
    \item Automatisation intégrale de la chaîne CI/CD (GitLab CI + Ansible + scripts Bash) avec mécanismes de rollback automatique.
    \item Intégration de la sécurité dès la conception (\emph{Shift Left Security}) : analyse statique SAST et détection de failles avec Snyk.
    \item Administration du stockage et des files de messages Cloud AWS (SQS, SNS, S3, Docker Registry).
    \item Configuration du système de traces et de surveillance applicative (Google Cloud Monitoring, Laravel Telescope).
\end{itemize}

\vspace{0.4em}

\textbf{\large Société Commerciale de Banque Cameroun (Attijariwafa Bank)} \hfill {\bfseries\small\color{slateGray} Avril 2022 -- Septembre 2022} \\
\emph{\color{primary}Architecte Infra \& DevSecOps (Projets Réglementaires)} \hfill {\small Cameroun}
\begin{itemize}[leftmargin=1.4em, itemsep=0.18em, topsep=0.15em]
    \item Déploiement et configuration de la passerelle de sécurité API KONG Gateway (Reverse Proxy, Load Balancing, limitation de débit, filtrage IP, protection DDoS).
    \item Déploiement d'un cluster Kubernetes avec Argo CD pour la livraison continue sur serveurs physiques HP ProLiant.
    \item Sécurisation des accès aux applications bancaires par authentification LDAP et gestion fine des habilitations RBAC.
\end{itemize}

% =========================================================================
% PROJETS OPEN SOURCE & PUBLICATIONS
% =========================================================================
\section{\faRocket\ Projets DevOps \& Publications Techniques}
\begin{itemize}[leftmargin=1.4em, itemsep=0.25em, topsep=0.15em]
    \item \textbf{webmethods-devops-portfolio :} Automatisation d'infrastructures d'intégration Software AG webMethods 10.x avec Terraform, Ansible et Docker.
    \item \textbf{mise-devtools :} Scripts d'automatisation et de provisionnement d'environnements de développement sous Linux et macOS.
    \item \textbf{Blog Technique Medium (\href{https://medium.com/@nkaurelien}{\texttt{medium.com/@nkaurelien}}) :} Rédaction d'articles spécialisés sur la sécurisation Nginx/SSL, le déploiement webMethods 10.15 et l'administration système Linux.
\end{itemize}

% =========================================================================
% FORMATIONS & DIPLÔMES
% =========================================================================
\section{\faGraduationCap\ Formations \& Diplômes}
\begin{itemize}[leftmargin=1.4em, itemsep=0.25em, topsep=0.15em]
    \item \textbf{2025} : \textbf{Formation webMethods 10.x} --- WILOW Paris (IS, Designer, Command Central, Administration).
    \item \textbf{2022 -- 2024} : \textbf{Master of Science (MSc) IA \& Management} / Expert en Ingénierie Informatique (RNCP 7) --- Ionis-STM France.
    \item \textbf{2014 -- 2017} : \textbf{Ingénieur des Travaux en Analyse \& Conception SI} --- Institut Africain de l'Informatique (IAI).
\end{itemize}

\vspace{0.8em}
\begin{center}
    \footnotesize\color{slateGray}
    CV édité sous \LaTeX\ $\cdot$ Retrouvez tous les projets et démos interactives sur \href{https://nkaurelien.kamitbrains.fr}{\textbf{nkaurelien.kamitbrains.fr}}
\end{center}

\end{document}
